\section{Kth Largest Element in an Array}
\label{sec:heap-kth-largest}

\problemstatement{Given an array $\textit{nums}$ of $n$ integers and an
integer $k$, find the $k$-th largest element in the array (the element
that would be at index $n-k$ if the array were sorted in increasing
order). Note this is the $k$-th largest \emph{value}, counting duplicates
by position, not the $k$-th distinct value.}

\begin{patternbox}
\emph{``Kth largest''} is the single most direct trigger for the
size-bounded heap technique introduced in
Section~\ref{sec:heap-intro}'s keyword list. The key (and slightly
counter-intuitive) design choice: to find the $k$-th \textbf{largest}
element efficiently, maintain a \textbf{min}-heap, not a max-heap -- the
min-heap's root is the \emph{worst} element among your current best-$k$
candidates, which is exactly the one you want instant access to evict.
\end{patternbox}

\subsection*{Understanding the problem carefully}

Sorting the whole array and reading off index $n-k$ costs
$O(n \log n)$ and uses the full order of all $n$ elements, even though we
only care about one specific rank. A heap lets us track only the $k$
largest elements seen so far, discarding everything else immediately.

\begin{ideabox}[Why a Min-Heap, not a Max-Heap]
Maintain a min-heap that never holds more than $k$ elements -- think of it
as ``the best $k$ candidates for being among the $k$ largest, with the
\emph{worst} of those $k$ candidates sitting at the root, ready to be
evicted.'' Process the array one element at a time: push the new element,
then if the heap's size exceeds $k$, pop the root (the smallest among the
current top-$k$ candidates, which is therefore the easiest one to safely
discard). After processing every element, the heap holds exactly the $k$
largest elements of the whole array, and its root -- the smallest of
those $k$ -- is precisely the $k$-th largest element overall.
\end{ideabox}

\subsection*{Why evicting the heap's root is always a safe discard}

At any point during the scan, the min-heap holds some $k$ candidates for
``currently in the top $k$.'' When a new element arrives and the heap
already has $k$ elements, we must decide: is the new element large enough
to belong in the top $k$, displacing the current weakest member? Pushing
the new element first and then popping the minimum handles both cases
uniformly without an explicit if-statement: if the new element is smaller
than everything already in the heap, it becomes the new minimum and is
immediately popped right back out, leaving the heap unchanged in effect;
if the new element is larger than the current minimum, that old minimum
(now provably not among the true top $k$, since at least $k$ other
elements, including the new one, are now known to be at least as large)
is correctly evicted instead. Either way, the heap always ends each step
holding the $k$ largest elements seen \emph{so far}, by induction.

\subsection*{Algorithm}

\begin{enumerate}
  \item Initialize an empty min-heap.
  \item For each element $x$ in $\textit{nums}$:
  \begin{enumerate}
    \item Push $x$ onto the heap.
    \item If the heap's size exceeds $k$: pop the minimum.
  \end{enumerate}
  \item Return the heap's current minimum (its root).
\end{enumerate}

\begin{algorithm}[H]
\caption{Kth Largest Element}
\begin{algorithmic}[1]
\Procedure{KthLargest}{\textit{nums}, $k$}
  \State $\textit{heap} \gets$ empty min-heap
  \For{each $x$ in \textit{nums}}
    \State push $x$ onto \textit{heap}
    \If{$|\textit{heap}| > k$}
      \State pop the minimum from \textit{heap}
    \EndIf
  \EndFor
  \State \Return top of \textit{heap}
\EndProcedure
\end{algorithmic}
\end{algorithm}

\subsection*{Dry run}

\dryrun{Take $\textit{nums} = [3,2,1,5,6,4]$, $k = 2$.}

\begin{itemize}
  \item Push $3$: heap $=\{3\}$. Size $1 \le 2$.
  \item Push $2$: heap $=\{2,3\}$. Size $2 \le 2$.
  \item Push $1$: heap $=\{1,2,3\}$. Size $3>2$: pop min ($1$). Heap
        $=\{2,3\}$.
  \item Push $5$: heap $=\{2,3,5\}$. Size $3>2$: pop min ($2$). Heap
        $=\{3,5\}$.
  \item Push $6$: heap $=\{3,5,6\}$. Size $3>2$: pop min ($3$). Heap
        $=\{5,6\}$.
  \item Push $4$: heap $=\{4,5,6\}$. Size $3>2$: pop min ($4$). Heap
        $=\{5,6\}$.
\end{itemize}

The final heap holds $\{5,6\}$, the two largest elements; its root
(minimum of these two) is $5$, the $2$nd largest element overall. Sorting
$[1,2,3,4,5,6]$ confirms index $n-k=4$ holds $5$.

\subsection*{C++ implementation}

\begin{lstlisting}
#include <bits/stdc++.h>
using namespace std;

int kthLargest(vector<int>& nums, int k) {
    priority_queue<int, vector<int>, greater<int>> minHeap; // min-heap

    for (int x : nums) {
        minHeap.push(x);
        if ((int)minHeap.size() > k) {
            minHeap.pop();
        }
    }
    return minHeap.top();
}

int main() {
    vector<int> nums = {3, 2, 1, 5, 6, 4};
    cout << "2nd largest: " << kthLargest(nums, 2) << endl;
    return 0;
}
\end{lstlisting}

\begin{complexitybox}
\textbf{Time Complexity.} Each of the $n$ elements triggers one push, and
at most one pop, each costing $O(\log k)$ since the heap never exceeds
size $k$, giving an overall time complexity of
\[
  O(n \log k),
\]
a clear improvement over $O(n \log n)$ when $k \ll n$.
\textbf{Space Complexity.} $O(k)$, since the heap never holds more than
$k$ elements.
\end{complexitybox}

\begin{takeawaybox}
The ``opposite polarity'' rule is worth memorizing precisely because it
feels backward at first: for \emph{largest}-$k$ problems, use a
\emph{min}-heap of size $k$ (so the easiest-to-discard candidate is always
at the root); for \emph{smallest}-$k$ problems
(Section~\ref{sec:heap-kth-smallest}), use a \emph{max}-heap of size $k$.
The heap's root is always the \emph{boundary} element of your current
best-$k$ set -- the one you're most willing to give up.
\end{takeawaybox}
